% To be inserted in the Methods section

\subsection{External Validation with Independent Mouse-Tracking and German Norm Data}

To test whether the fitted response-competition parameter generalized beyond the KH2017 semantic categorization dataset, we analyzed the public mouse-tracking archive from Koenig-Robert et al. (2024; OSF: \url{https://osf.io/q3hbp/}). The downloaded files were \texttt{data\_face.zip}, \texttt{data\_animal.zip}, and the authors' MATLAB preprocessing and analysis scripts. Downloads were obtained on 18 June 2026; file hashes and preprocessing counts are reported in \texttt{results/external\_koenig\_robert\_preprocessing\_report.md}. The face/object study included true faces, ordinary objects, and face-like objects that were to be categorized as objects. The animal/object study included true animals, ordinary objects, and animal-like objects that were to be categorized as objects. Extra balancing stimuli distributed with the experiment were not included in the primary analyses.

Raw Pavlovia/jsPsych CSV files were parsed directly from the OSF archives. We retained complete participant files, applied the public preprocessing script's response-side counterbalancing, excluded invalid trials and participants using the authors' validity and accuracy criteria, excluded incorrect trials for model fitting, and canonicalized each valid trajectory to the KH2017 geometry: start at lower center, selected target at upper right, and competitor at upper left. The action model used the same numerical settings as the KH2017 analysis: 51 time samples and a $\rho$ grid from 0 to 2 in steps of 0.05.

The primary external tests compared fitted $\rho$ for lookalike objects versus ordinary objects within each study. Because the downloaded OSF archive did not include scalar face-likeness or animal-likeness ratings, lookalike-object status was treated as a prespecified binary ambiguity proxy rather than as an independent continuous rating. We used stimulus-level mean $\rho$, participant-clustered and stimulus-clustered bootstrap intervals, stimulus-label permutation controls, leave-one-stimulus-out held-out stochastic path likelihood, and pooled-object null simulations with no stimulus-level ambiguity predictor. Flexible raw-trajectory baselines were retained as descriptive benchmarks, and lower raw NLL by such baselines was not treated as evidence against the interpretability of $\rho$.

We also obtained SWOW-DE 2025 from Small World of Words (\url{https://smallworldofwords.org/en/project/research}) on 18 June 2026, using the recommended R55 release and the official PPMI--SVD embedding bundle. KH2017 German item and category labels were canonicalized before matching (e.g., \emph{Saeugetier} to \emph{Säugetier}, \emph{Chamaeleon} to \emph{Chamäleon}, \emph{Seeloewe} to \emph{Seelöwe}, and \emph{Loewe} to \emph{Löwe}). Direct cue-to-response, reverse category-to-item, symmetric association, undirected graph-proximity, and official embedding-cosine margins were computed. Before inspecting the $\rho$ association, the main SWOW metric was defined as direct association if direct coverage reached at least 12 of 19 items; otherwise, network or embedding proximity was used when it covered at least 16 items.

Finally, Schröder et al.'s German semantic norms were audited as lexical-semantic controls. These norms include broad typicality, familiarity, age of acquisition, frequency, and word length measures, but they do not provide exact KH2017 target--competitor category typicality contrasts. They were therefore used only as coverage and control variables, not as replacements for the semantic-margin predictor.
